This paper documents a implemented pipeline that combines case retrieval, adaptive capability weighting, bottleneck detection, reflection, structured simulation evaluation, and standards-oriented exporters for simulated multi-domain operation planning. Using a single canonical evidence scope, the paper reports that the Full Memento--Hope--Reflection configuration outperforms the pure LLM baseline in the selected ablation study and in four of five scenarios of the selected cross-scenario transfer suite (RTX 4090 D evidence).

Beyond the four contributions enumerated in the Introduction, several enhancements strengthen the current work. First, a \texttt{BottleneckDetector} module replaces the static bottleneck-boost vector with a dynamic, history-aware diagnosis that identifies the current bottleneck stage and allocates capability emphasis proportionally to measured severity. The detector flags \texttt{breach} as the persistent bottleneck in 9 of 10 iterations, and the weight-competition analysis reveals that breach and sustain share the \texttt{protection} capability dimension, creating a structural tension that a scalarized reward cannot fully resolve. Second, a memory-conditioned Hope coupling lets retrieved case metrics bias the fast-weight initialization, so the pipeline's two auxiliary modules share information through a common capability representation. Third, multi-seed replication across five seeds confirms the Full pipeline's gain is statistically significant under default HOPE parameters (paired $t(4)=3.75$, $p=0.02$, $d=1.68$; overall-effectiveness gain $p<0.01$). Fourth, external baselines, a three-backend comparison, and a Reflection-only single-scenario ablation provide converging evidence for the pipeline's effectiveness and clarify that standalone reflection degrades plan quality under the fixed-seed setting.

Just as importantly, the paper makes its current boundaries explicit. It does not claim formal proofs, external benchmark leadership, or evidence beyond what the current evidence can support. The multi-seed evidence, baseline comparisons, three-backend validation, and Reflection-only replication reported here substantially narrow the gap between the original framework and the statistical rigor expected of a journal submission. Several limitations remain: (1) the cross-scenario CIs reflect only within-seed Monte Carlo variance; (2) the per-component ablation (Memento-only, Hope-only) is multi-seed only for seeds 7 and 42; and (3) the case-bank size ablation (Section~4.10) is currently reported for a single seed. These limitations define directions for future work and do not invalidate the current contribution.